\documentclass[11pt]{article}
\usepackage[T1]{fontenc}
\usepackage{amsmath,amssymb,geometry,hyperref}
\geometry{margin=1in}
\newcommand{\lean}[1]{\texttt{\detokenize{#1}}}
\newcommand{\K}{\mathsf K}
\title{Weak Grothendieck payloads: actual operations and the horizontal coherence boundary}
\date{19 November 2026}
\begin{document}\maketitle
\paragraph{Verdict and scope.} The package constructs the endo-2-functor, its pseudonatural unit and multiplication, and the invertible components of all three monad constraints on the actual composites. It does not bundle or assert a completed pseudomonad. Two different equations previously described as the pentagon can now be distinguished by their actual Lean types. The canonical native schedule equality lifts to an equality of modifications in the presheaf variable, and is reflected there. An equality of horizontal pastings, whose short side uses the naturality of multiplication at multiplication, is a different statement. The remaining comparison is localized below. The paper's \lean{eq:grothmonad} is at issue, not its subsequent \lean{prop:groth} on representable Kleisli restriction. No representability or Kleisli hypothesis or conclusion is used.
\paragraph{Objects fixed; arrows not strictified.} For a fixed ordinary base $B$ and split graded endpoint datum $P,S$, the objects are ordinary functors $B^{\mathrm{op}}\to\mathbf{Cat}$. Hom categories are the full hom categories of the corresponding mathlib pseudofunctors: arrows are strong transformations; 2-arrows are arbitrary modifications. \lean{AmbientBicategory.Obj} is induced along this inclusion. Its unit and association isomorphisms are equality transports; this says nothing about a payload map's reindexing isomorphisms being casts. For the finite block actions, $\K_0=\mathrm{id}$ and $\K_n;\K_1=\K_{n+1}$ as entire strict pseudofunctors, including action on noninvertible modifications. Thus $\K=\K_1$, and the direct two-block action has the genuine composite source $\K^2$. Length zero is not inferred from positive lengths.
\paragraph{Comparison mechanism.} Encode a strong map by its functor on co-Grothendieck totals over $B$. Decode a cartesian functor by choosing cartesian comparisons, not by pretending the functor preserves chosen lifts on the nose. \lean{AmbientRoundTrip} proves both round trips on arrows including moving-base comparators; \lean{AmbientVerticalDescent} handles noninvertible vertical transformations. \lean{AmbientWhiskering} proves both horizontal whisker comparisons by faithful grade-word and terminal-edge observation. \lean{AmbientMateCalculus.factor} characterizes the weak weighted square by composition with the directed terminal cartesian lift. The mapped and chosen lifts need only be cartesian over the same arrow. Two projections and cancellation prove modification naturality, identity and composition, not just objectwise isomorphisms. No endpoint-leg inverse is used.
\paragraph{Unit, multiplication and modifications.} \lean{AmbientPseudonaturality.operation} gives a strong transformation $\K_i\Rightarrow\K_j$ for every native cartesian operation $r:i\to j$. \lean{AmbientMonadMaps.eta} has source the identity; \lean{AmbientMonadMaps.mu} has source $\K;\K$. \lean{AmbientOperationBoundaries.nativeMultiplication} compares descended multiplication with the previously named strong multiplication, including its exterior squares. \lean{AmbientNativeModification.iso} sends a native cell to an invertible modification between the payload-pseudonatural operations, not only a family of modifications over the base. \lean{AmbientScheduleRealization.realize} pastes the actual ambient modifications; \lean{faithfulPaths} reflects equality by decorating an arbitrary grade row at its outgoing endpoint using the constant observer. Other observers with empty fibers are not presumed faithful.
\paragraph{The horizontal equations.} Write $M_X=\mu_X$, $E_X=\eta_X$, in travel order. The invertible cells of \lean{AmbientHorizontalFrames} are
\[
 L_X:E_{\K X};M_X\cong1,\qquad
 R_X:\K E_X;M_X\cong1,\qquad
 A_X:M_{\K X};M_X\cong\K M_X;M_X .
\]
They are constructed by cartesian descent and identified component boundaries, not postulated. Set $a=M_{\K^2X}$, $b=\K M_{\K X}$, $c=\K^2M_X$, $d=M_{\K X}$, $e=\K M_X$, $t=M_X$. The horizontal pentagon compares
\[
 (a d)t\longrightarrow (b d)t\longrightarrow (b e)t\longrightarrow (c e)t,\qquad
 (a d)t\longrightarrow (a e)t\longrightarrow (c d)t\longrightarrow (c e)t .
\]
The middle short step uses the inverse naturality comparison of $M$ at $M_X$, whiskered by $t$. It is not specified to be a cast. The long steps use $A_{\K X}$, $A_X$ and $\K(A_X)$ respectively. The triangle starts at $(\K E_{\K X};M_{\K X});M_X$ and compares the switch by $A_X$ then $\K(L_X)$ with $R_{\K X}$ whiskered by $M_X$. Actual unitors, associators and structural map comparisons are retained in \lean{AmbientHorizontalDiagrams}; the strings above suppress them for display.
\paragraph{Not a missing axiom on the target.} The three naturality squares required of $L,R,A$ are \lean{AmbientHorizontalNaturality.Left}, \lean{Right}, and \lean{Assoc}, with the source and target comparisons formed from the actual unit/multiplication comparisons by whiskering. \lean{AmbientSquareBridge.viaNative} suffices once these horizontal boundary comparisons are identified with the chosen weak mates of the associated serial native operations, with the proved component frames on both sides. Native constraint naturality then supplies the interior, without hypothesizing target naturality. Component frames by themselves do not imply this square identification. Likewise canonical native schedule equality cannot be substituted directly into either horizontal diagram.
\paragraph{Discriminating probes.} \lean{DirectedModification.lean} constructs a modification $0\le1$ between constant weak transformations on a constant poset payload. A reverse cell at a nonempty row of any finite arity would entail $1\le0$. \lean{AmbientUnitMate.lean} chooses a split endpoint extending to $0\in\mathbb N$ and directed discharge $0\to i$. A payload with fiber $B C_2$ has identity component functors but a signed weak naturality isomorphism across $0\to1$. The actual unit operation's comparison has terminal bit one at the selected row; an equality transport has bit zero. Thus replacing weak squares by casts fails on the operation to be identified, not just unrelated abstract data. The separate \lean{HorizontalDemand.lean} probe reuses the pinned native signed model: the triangle now holds and the pentagon fails for the vertically realized presheaf modifications, by faithful reflection, not just for a bare grade test. A failure to reduce by definitional equality is not used as evidence.
\paragraph{One local test before the final multiplication.} \lean{AmbientExchangeDemand.interchange} is the inverse of the actual multiplication comparison on $M_X$. Its source is $a;e$, its target is $c;d$. \lean{AmbientExchangeDemand.Law} asks that, on a decorated four-grade row before final multiplication, this particular arrow read the pinned canonical disjoint exchange of the two native contractions. This is equality up to the actual source and target transports on each row, not an arbitrary isomorphism chosen to make the pentagon commute. The complete source and destination \emph{functors} already match: in \lean{AmbientDisjointFrames}, \lean{disjointLeft} identifies the encoding of $a;e$ with the native left disjoint route and \lean{disjointRight} identifies the encoding of $c;d$ with the right one. These are equalities on objects and moving arrows, not just agreement on word lengths. The unproved primitive at that vertex is the image of the naturality \emph{cell} between those functors. Its post-whiskering implication is proved in \lean{whiskered}. In the old moving-color model the native exchange crosses unequal shade objects, even though the colors are isomorphic. \lean{ExchangeFalsifier.lean} rules out every cast as an answer to this local comparison test in that model. The independent check \lean{MovingAmbientExchange.lean} goes further: decorate the four-grade row by a point of the constant observer, then evaluate the \emph{actual} two sides $a;e$ and $c;d$. The \lean{different} declaration proves their resulting objects unequal, using the left-white/right-black grade calculation and both full endpoint functor identifications. Consequently no arrow between these objects can be an equality transport (the file separately proves the general criterion \lean{noCastAt}). This does not decide which nontrivial isomorphism occurs: that is precisely the residual cell identification.
\paragraph{Constructive reduction and first discharged boundaries.} \lean{AmbientGeneratorCriterion.Matches} is an inductive certificate assembled one fixed native generator and one actual arrow at a time. Its soundness and transfer lemmas turn two such local identifications plus the native law into equality of actual pastes. They cannot infer a diagram from equality of its endpoints and do not allow adjustment of the canonical finish. \lean{AmbientHorizontalRead.pre} and \lean{post} discharge ordinary serial whiskers, subject to the explicit, checkable chart equality for the prefix/suffix arrow. \lean{AmbientLetterRead} reads the three component constraints back to their prescribed native cells. \lean{AmbientIdentityOperation.comparison} proves that the weighted identity operation's weak comparison is an equality transport for \emph{every} weak arrow, using the direction-sensitive mate calculation with identity legs; this statement emphatically does not extend to the native unit or tensor.
\paragraph{Route for independent rejection or continuation.} Build against the pinned package with no project oleans, then run the signed and directed probes. At the mathematical boundary, check \lean{AmbientHorizontalNaturality} and \lean{AmbientHorizontalDiagrams} first, so that the target composites and orientations are fixed independently of native seams. Read \lean{AmbientDisjointFrames} and \lean{AmbientExchangeDemand.Law} before looking at the native proof of \lean{ContextadDisjointCut.disjoint}. A possible next proof tests the two projections of the actual decoded comparison against the defining cartesian factorization of that native cell. Endpoint equality alone, an arbitrary invertible filler between those endpoints, or a pentagon satisfied only after applying the last multiplication is weaker and cannot establish this comparison. For the unitors, use the common weighted-identity comparison already checked above; on the composite side the serial comparison must be reconciled with the inner/outer whiskered unit comparisons. Independence of triangle and pentagon in the signed model prevents deleting the source pentagon assumption. No claim that coherent native input is insufficient, nor an impossibility theorem for the remaining bridges, follows.
\end{document}
